% =====================================================================
% Response to supervisor's marked corrections on Paper A and Paper C.
%
% Standalone: compiles with a bare `pdflatex response_to_supervisor.tex`
% (twice, for the reference to the item count). No bibliography, no
% external figures -- it must survive being mailed on its own.
% =====================================================================
\documentclass[11pt,a4paper]{article}

\usepackage[T1]{fontenc}
\usepackage[utf8]{inputenc}
% lmodern is absent from this toolchain; default CM at T1 is fine
\usepackage[expansion=false]{microtype}
\usepackage[margin=2.4cm]{geometry}
\usepackage{amsmath,amssymb}
\usepackage{booktabs}
\usepackage{enumitem}
\usepackage{xcolor}
\usepackage{titlesec}
\usepackage[hidelinks]{hyperref}

\definecolor{qgrey}{HTML}{4A4A4A}
\definecolor{rule}{HTML}{C9C9C4}
\definecolor{tag}{HTML}{8A4B2A}

\titleformat{\section}{\large\bfseries}{\thesection}{0.7em}{}
\titlespacing*{\section}{0pt}{1.6em}{0.7em}

\setlength{\parskip}{0.35em}
\setlength{\parindent}{0pt}

% One response item. #1 = tag, #2 = location in the marked file,
% #3 = the supervisor's query, #4 = the response, #5 = where it landed.
\newcommand{\thinrule}{\textcolor{rule}{\rule{\linewidth}{0.4pt}}}
\newenvironment{qitem}[3]{%
  \par\medskip\thinrule\par\smallskip
  \noindent\textbf{\textcolor{tag}{#1}}\quad\textcolor{qgrey}{\small #2}\par\smallskip
  \noindent\textcolor{qgrey}{\small\itshape #3}\par\smallskip
}{\par}
\newcommand{\where}[1]{\par\smallskip{\small\textcolor{tag}{$\hookrightarrow$}\ \textit{#1}}\par}

\begin{document}

\begin{center}
{\Large\bfseries Response to Reviewer Comments}\\[0.35em]
{\large Papers A and C, consolidated into the merged submission}\\[0.9em]
\begin{tabular}{@{}ll@{}}
\textbf{From} & Roger Nick Anaedevha\\
\textbf{Re} & 20 marked queries on \texttt{paperA.tex} (6) and \texttt{paperC.tex} (14)\\
\textbf{Outcome} & All 20 addressed in the merged paper, target NDSS 2027\\
\end{tabular}
\end{center}

\vspace{0.6em}

\section*{Summary}

Thank you for the detailed reading. Every marked query is addressed below, with
the location in the merged paper where the change landed. Four points frame the
rest.

\textbf{The two papers are now one.} Your comment that reviewers receive only
one paper and cannot know what is in the other was decisive. The benchmark and
the theorem are now a single submission whose thesis is the end-to-end
framework; the benchmark appears as the instrument that makes the composed
guarantee measurable, not as a dataset release. The phrase ``companion paper''
no longer occurs anywhere in the manuscript.

\textbf{Three queries changed the technical content, not just the presentation.}
The temporal MCR predicate (C9), the alignment error bound (C11), and the ROC
generalisation (C13) each required new derivations. They are the strongest
additions in the merged paper and we are grateful for them.

\textbf{One query produced a negative result we have reported rather than
hidden.} The provenance breakdown you asked for (C14) shows that the number of
released pairings in our strongest evidence class is \emph{zero}. We state this
in the abstract and in the limitations section.

\textbf{On the count.} The marked files contain 20 \verb|\fc{...}| annotations:
6 in \texttt{paperA.tex} and 14 in \texttt{paperC.tex}. Two pairs sit on a
single line each (C1/C2 and C6/C7), which we believe accounts for the count of
15 you may have in mind. All 20 are answered; none has been merged away or
skipped.

% =====================================================================
\section*{Paper C — benchmark and harness (14 queries)}
% =====================================================================

\begin{qitem}{C1}{abstract, line 54}{``In general, certain phrases are too
conversational for a scientific paper.''}
Accepted throughout, not only at the marked point. We removed the rhetorical
register wholesale: figurative framing (``Achilles heel''), first-person
narration of the research process, and sentences whose work was emphasis rather
than content. Where a sentence existed to signal enthusiasm it was deleted;
where it carried a claim, the claim was restated declaratively. The abstract was
rewritten to state the result rather than to introduce it.
\where{Merged paper, throughout; the abstract in full.}
\end{qitem}

\begin{qitem}{C2}{abstract, line 55}{``If we combine the papers the central
thesis must be the end-to-end framework. The benchmark should not be presented
merely as a dataset release.''}
Accepted, and it reorganised the paper. The thesis is now the composition
$\theta \mapsto \Delta(\theta) \mapsto \delta(\theta) \mapsto \mathrm{MCR}$, and
the six-source corpus is introduced as the measurement instrument that closes
the interface between the two layers---specifically, the only way to obtain the
delay-to-position rate $\gamma$ empirically rather than assume it. The title now
names the framework: \textit{From Wire to Flight: An End-to-End Framework
Composing Network Intrusion Detection with Certified Mission Safety for
Autonomous UAVs}. The dataset contribution is stated as an enabling artifact.
\where{Title, abstract, Introduction; the contribution list is reordered so the framework is first.}
\end{qitem}

\begin{qitem}{C3}{§Introduction, line 157}{``I would prefer to merge the two
papers. We should never mention `companion paper' in a scientific paper---if we
send two separate papers the reviewers will only have one, they cannot know
what's in the other one.''}
Accepted without reservation; this was the correct call and we had
underestimated it. The scoping sentence and every forward reference to a second
paper are gone. The merged manuscript is self-contained: the theorem is proved
in it, and the benchmark that instantiates the theorem is described in it.
\where{Zero occurrences of ``companion paper'' remain in the source.}
\end{qitem}

\begin{qitem}{C4}{§Background, line 164}{``The introduction to this section must
be rewritten.''}
Rewritten, and relocated. The section opened by announcing its own structure,
which told the reader nothing. It now opens with the substantive claim the
section exists to establish---that the attacks which matter in this field are
network events whose damage is an autonomy outcome, so neither layer's
literature can measure them alone. To meet the NDSS 13-page body limit without
cutting content, the tutorial material moved to an appendix and only the load-bearing
definitions remain in the body.
\where{§II opening paragraph; expanded background in Appendix~A.}
\end{qitem}

\begin{qitem}{C5}{§``Anatomy of an autonomous UAV under attack'', line 169}{``Rename
this section.''}
Renamed. The heading was descriptive of a narrative rather than of content. It
is now stated as the system and threat model it actually presents.
\where{§II-A.}
\end{qitem}

\begin{qitem}{C6}{Table of acronyms, line 189}{``I would prefer to remove this
table. We define the acronyms in the text and then we use them.''}
Removed. The table is deleted and each acronym is expanded at first use in
running text.
\where{No \texttt{tab:acronyms} remains in the source.}
\end{qitem}

\begin{qitem}{C7}{same line}{``Note also that at the moment you (almost
everywhere) are doing both things---defining the acronym in text and having the
table for it.''}
Correct, and the duplication is what made the table redundant rather than merely
long. We audited every acronym: each is now expanded exactly once, at first use,
and used bare thereafter. Acronyms appearing only once are spelled out and the
abbreviation dropped.
\where{Throughout.}
\end{qitem}

\begin{qitem}{C8}{Fig.~``attack surfaces'', line 269}{``There is a weird arrow on
the top-left square. Is it useful? I don't understand where it points to.''}
Not useful---it was a drafting artifact with no referent, and we should have
caught it. Removed. The figure now carries exactly one downward arrow, the
network-to-autonomy coupling that is the paper's subject, and the caption names
it explicitly so the single remaining arrow cannot be misread.
\where{Fig.~1 and its caption.}
\end{qitem}

\begin{qitem}{C9}{§``Mission Completion Rate'', line 274}{``The paper adopts a
purely spatial definition of MCR. But a drone that safely finishes its route yet
arrives 30 minutes late due to an evasive time-delay attack has still failed a
time-critical mission. Augment the MCR definition with a temporal bounding
constraint, or explicitly define it as `Spatial MCR'.''}
Accepted, and this changed the theorem rather than only the definition. The
observation is sharp: a delay attack that our own threat model treats as the
primary adversary was, under the old definition, capable of succeeding
completely while the metric recorded success. MCR is now a conjunction of a
spatial and a temporal predicate,
\[
\mathrm{MCR}=\Pr\Big[\underbrace{\sup_{t\in[0,T]}\lVert x(t)-x_{\mathrm{nom}}(t)\rVert\le m}_{\text{spatial}}
\;\wedge\;
\underbrace{t_{\mathrm{arr}}\le(1+\kappa)T}_{\text{temporal}}\Big],
\]
with $\kappa$ the mission's schedule tolerance. The composition theorem now
discharges both: the spatial predicate via the Grönwall tube as before, and the
temporal predicate via the condition $\Delta(\theta)\le\kappa T$ on the same
residual delay budget the detector already bounds. The two predicates are
therefore certified by one quantity, which is a stronger result than the
original and follows directly from your comment.
\where{§III-C (definition); Theorem 1 and its proof; the temporal condition is discharged in the same lemma as the spatial one.}
\end{qitem}

\begin{qitem}{C10}{§Schema, line 412}{``Eq.''}
Read as a request to display and number the pairing tuple rather than leave it
inline. Done: the \textsc{CrossLayerPairing} record is now a numbered equation
and is referenced by number where the certificate consumes it.
\where{Eq.~(3) and its two later references.}
\end{qitem}

\begin{qitem}{C11}{§Schema design, line 425}{``You state that time is always
relative to each source's own capture start. Add a paragraph mathematically
defining how you align relative network time to relative telemetry time for the
\texttt{synthetic\_alignment} pairings. Detail the synchronization error bounds,
or explain the heuristic that guarantees synthetic alignment does not
artificially inflate or deflate $\Delta(\theta)$.''}
Accepted; a derivation has been added. We anchor attack onsets under an affine
map between the two relative clocks, decompose the residual into three terms
(onset-detection error in each source, plus drift over the window), and bound
the total at $e_{\mathrm{align}}\le0.53$\,s for every released pairing, with a
rejection rule that discards any pairing exceeding it.

One point deserves emphasis because it answers the sharper half of your
question. Misalignment cannot inflate or deflate $\Delta(\theta)$, because
$\Delta(\theta)$ is computed \emph{entirely within the network source} from
per-hop residual delays; the alignment map is used only to attach that value to
an autonomy window. A bad alignment can therefore pair a real network
measurement with the wrong flight segment---which we bound and report---but
cannot corrupt the quantity the certificate consumes. We now state this
explicitly rather than leaving the reader to verify it.
\where{§IV-C with the bound; full map, three error terms and rejection rule in Appendix~B.}
\end{qitem}

\begin{qitem}{C12}{§``What each adapter does'', line 542}{``Rename the
section.''}
Renamed to describe the mapping the section actually specifies, from
heterogeneous source records to the unified schema, rather than the activity of
writing adapters.
\where{§IV-D.}
\end{qitem}

\begin{qitem}{C13}{§Harness, line 618}{``You build the harness entirely around a
DATAMUt-style detector with a deterministic threshold $\epsilon$, while others
rely on ML classifiers. Add a discussion or small empirical subsection showing
how your schema and harness map to ML-based ROC curves---how a classifier's
confidence cutoff $P_{\mathrm{th}}$ maps to the residual undetected budget. This
proves your benchmark is a framework for any IDS, not just DATAMUt.''}
Accepted; this materially widens the paper's claim and we have added the
subsection. For a scoring detector with score function $g$ and cutoff
$P_{\mathrm{th}}$, an evading attacker must satisfy $g(\mu)<P_{\mathrm{th}}$,
hence
\[
\mu \;<\; g^{-1}(P_{\mathrm{th}}) \;=:\; \theta_{\mathrm{eff}}(P_{\mathrm{th}}),
\]
and the residual-budget lemma applies verbatim with $\theta$ replaced by the
effective operating point $\theta_{\mathrm{eff}}$. The threshold detector is the
special case $g=\mathrm{id}$, for which $\theta_{\mathrm{eff}}=\varepsilon$. We
state the two properties of $g$ the composition needs---monotonicity, so that a
stricter cutoff is a smaller budget, and a measurable inverse on the operating
range---and note that any detector emitting a per-hop score satisfies them.

The consequence is exactly the one you identified: the certified floor is a
function of a detector's ROC operating point, not of DATAMUt. Any IDS that
reports an ROC curve can be substituted without touching the theorem.
\where{§V-B, with Eq.~(7); the DATAMUt instantiation becomes an example rather than the definition.}
\end{qitem}

\begin{qitem}{C14}{§Evidence classes, line 830}{``Add a summary table or clear
breakdown showing the exact distribution of provenance tags across the 93,600
flights. Reviewers will want to know exactly what percentage relies on synthetic
alignment versus actual physical measurements.''}
Added, and it is the most uncomfortable table in the paper. We built a
provenance ledger over the whole corpus and report the distribution: of 431,773
ingested events, 200,610 (46.5\%) originate from real hardware captures.

The breakdown also surfaces a fact we had not previously stated plainly, and we
have chosen to report it prominently rather than let a reviewer discover it: the
number of released pairings tagged \texttt{measured\_same\_platform}---our
strongest evidence class, requiring both layers observed on one platform---is
\textbf{zero}. The single source that observes both lacks machine-readable
attack intervals, so its pairings cannot yet be released in that class. The
$\gamma$ calibration underlying the headline result is still a real-flight
measurement, from three PX4 flights, and we now label it as a calibration sample
rather than a distribution. This is stated in the abstract, tabulated in §VI,
and listed first among the limitations.
\where{Table~IV (provenance distribution), abstract, §VIII Limitations.}
\end{qitem}

% =====================================================================
\section*{Paper A — composition theorem (6 queries)}
% =====================================================================

\begin{qitem}{A1}{title, line 37}{``We should specify the attack or remove the
keyword from the title.''}
Removed. ``under Attack'' promised a specificity the paper does not deliver, since
the theorem is stated for a class of network-induced perturbations rather than a
named attack. The merged title names the mechanism instead: \textit{From Wire to
Flight: An End-to-End Framework Composing Network Intrusion Detection with
Certified Mission Safety for Autonomous UAVs}.
\where{Title.}
\end{qitem}

\begin{qitem}{A2}{§$\gamma$ calibration, line 537}{``Explicitly define the
temporal validity window for $\gamma$: `We model this degradation as locally
linear for $\Delta(\theta)\le[X]$ seconds; beyond this threshold, uncorrected IMU
drift bounds become quadratic, requiring a higher-order $\gamma$ function.'\,''}
Accepted, with the constant measured rather than asserted. We state
$\tau_{\mathrm{lin}}\approx2$\,s, taken from the innovation-gating horizon
observed in the calibration flights: while GNSS is late rather than absent, the
EKF continues correcting against its last accepted fix and error grows
approximately linearly in staleness; once the outage exceeds that horizon the
filter propagates on inertial data alone and uncorrected accelerometer bias and
gyroscope drift make the error grow quadratically in time.

This is favourable to the paper's claim and we note why: the certified operating
window of interest is sub-second, well inside $\tau_{\mathrm{lin}}$, so the
linear model is used only where it is valid. We state the quadratic form for
completeness and flag the higher-order $\gamma$ as future work rather than
extrapolating the linear fit past its range.
\where{§III-D, ``Temporal validity of the linear model''.}
\end{qitem}

\begin{qitem}{A3}{Theorem 1, line 548}{``The Grönwall trajectory tube relies
heavily on the navigation drift $F$ being $L$-Lipschitz. However, this assumes
the drone's physical dynamics remain in a smooth, linearizable regime. Maybe we
should note this in the paper.''}
Noted, and given the standing of a hypothesis rather than a technicality. A
remark now follows the theorem stating that $L$-Lipschitzness holds \emph{over
the operating region} and that the qualifier carries physical content: attitudes
away from gimbal-lock singularities, airspeeds below stall or the onset of blade
stall, actuators off their saturation limits, and no contact dynamics. We state
where these hold and where they do not, and are explicit that an aggressive
evasive manoeuvre or actuator saturation exits the certified regime---so the
guarantee lapses rather than degrades.

Related, and prompted by the same scrutiny: we re-measured $L$ in the operating
region and obtained $1.18$, against the global value of $1.01$ used earlier. We
now report the local, larger constant throughout, which tightens the certified
radius to $0.153$. The headline result at a 10\,m corridor is unaffected; a
conclusion at a 2\,m corridor did change, and we state that in the limitations
rather than only reporting the margin that survives.
\where{Remark ``On the Lipschitz hypothesis'' after Theorem 1; constants table; §VIII.}
\end{qitem}

\begin{qitem}{A4}{§Discussion, line 629}{``Remember to remove the conversational
tone.''}
Done, and treated as a global pass rather than a local edit, together with C1.
The marked passage (``Achilles heel'') and its neighbours are rewritten as
declarative statements of the limitation they describe.
\where{Throughout; no figurative characterisations of results remain.}
\end{qitem}

\begin{qitem}{A5}{§Four-certificate architecture, line 676}{``The core theorem
solves a time-delay attack, while MWU is suddenly introduced for an `adaptive
online jammer' and PAC-Bayes awaits pending data. Write a unifying paragraph
explaining exactly how they operate simultaneously on the same flight---otherwise
severely trim the PAC-Bayes and MWU sections. Better one deeply proven,
empirically verified theorem than a survey of four loosely connected bounds.''}
This was the most consequential comment on Paper A and we accepted its framing
before choosing between the two remedies. We took the first: a paragraph now
holds one flight fixed---a ten-minute delivery sortie, flown by a model trained
on last season's mission mix, through a corridor where a compromised relay
delays peer corrections and a ground jammer varies transmit power---and shows
that four claims are in force simultaneously about four different objects: the
Grönwall tube bounds the trajectory, randomized smoothing bounds the classifier
under input perturbation, PAC-Bayes bounds generalisation from the training mix
to this sortie, and MWU bounds regret against the adaptive jammer. They are
concurrent, not alternative.

We also acted on the trimming half of your comment. PAC-Bayes and MWU are
reduced to what the unifying argument requires; the paper's weight sits on the
Grönwall result, which is the one proved and empirically verified. Where
PAC-Bayes awaits a numeric $\mathrm{KL}$ term we flag it as pending in the
constants table rather than substituting a plausible value.
\where{§VI, ``How the four operate on a single flight''; Table~III marks the pending constant.}
\end{qitem}

\begin{qitem}{A6}{§Instantiation, line 772}{``To make the evaluation realistic we
could cite actual aviation or drone regulatory standards (e.g.\ FAA requirements
for urban delivery drones or U-Space tracking tolerances) to justify why 5\,m or
10\,m are the correct thresholds. Do you think we can do that?''}
Yes, and it is done. The corridor margins were previously round numbers with no
external warrant. We now ground them in the JARUS Specific Operations Risk
Assessment (SORA), the reference methodology adopted by EASA and other
authorities for the \emph{specific} category, under which the operational volume
comprises the flight geography plus a contingency volume whose minimum lateral
extent is a sum of independent error terms. Using the standard's own default
values---3\,m GNSS position accuracy, 3\,m position-holding error, and 1\,m map
error---gives 7\,m before any reaction-distance or manoeuvre term.

The evaluated margins are now positioned against that figure rather than chosen:
2\,m is tighter than any SORA-compliant buffer, 5\,m remains below the default
contingency sum, and 10\,m is the first value clearing it with margin. The
certified window becomes comparable to a regulatory quantity instead of to a
modelling choice, which is what your question was asking for.
\where{§VII-A with the SORA derivation; margin table; the headline figure is annotated with the 7\,m line.}
\end{qitem}

% =====================================================================
\section*{Two items we flag rather than close}
% =====================================================================

Neither was raised in your comments; both follow from acting on them, and we
would rather surface them now than have them found at review.

\textbf{The 2\,m conclusion changed (from A3).} Re-measuring $L$ locally at
$1.18$ means $\theta=0.25$\,s is certified at corridor margins $m\ge5$\,m but
\emph{not} at $m=2$\,m, where the tube reaches $2.22$\,m. The earlier draft
claimed otherwise on the strength of the global constant. The headline 10\,m
result is unaffected.

\textbf{The strongest evidence class is empty (from C14).} As above: zero
released \texttt{measured\_same\_platform} pairings. We judged that reporting
this openly is more defensible than presenting the corpus in a way that does not
invite the question, and it is now stated in the abstract.

\vspace{1.2em}
\thinrule
\vspace{0.5em}

{\small All changes are in the merged manuscript. Every number quoted here is
reproducible from the committed artifact; the provenance ledger that produced
the C14 distribution is \texttt{experiments/provenance\_ledger.py}, and its
output is \texttt{results/provenance\_distribution.csv}.}

\end{document}
